% !TeX root = ../../Book1.tex
\section{测试函数}
\label{b1:sec:2001}

对光滑函数作分部积分，可以把一个函数的导数移到另一个函数上。若第二个函数在边界附近为零，边界项也随之消失。于是，即使第一个函数不可微，积分的另一边仍可能有意义。为了沿着这条思路推广导数，先要选择一类既能任意次求导、又能把计算限制在开集内部的函数。

\ProseParagraphBreak

本节固定非空开集 \(\Omega\subset\mathbb R^d\)，标量域为 \(\mathbb K=\mathbb R\) 或 \(\mathbb C\)。记号 \(K\Subset\Omega\) 表示 \(K\) 是包含于 \(\Omega\) 的紧集。本节的拓扑构造只使用半范数、紧集穷竭和光滑截断，不预借分布或 Sobolev 空间的性质。
% 来源边界：Hunter2001Applied第11章主要使用Schwartz空间；作者公开稿印刷294页明确区分D'与S'。
% 本节DK完备性、单位分解和D的局部凸拓扑由下文完整构造，不把该书温和分布版本冒称为这里的证明。

\subsection{紧支集光滑函数}
\label{b1:subsec:200101}

对多重指标 \(\alpha=(\alpha_1,\ldots,\alpha_d)\in\mathbb N_0^d\)，约定
\[
 |\alpha|\coloneqq\sum_{i=1}^d\alpha_i,\quad
 \alpha!\coloneqq\prod_{i=1}^d\alpha_i!,\quad
 D^\alpha\coloneqq\partial_1^{\alpha_1}\cdots\partial_d^{\alpha_d}.
\]
不等式 \(\beta\leq\alpha\) 指每个坐标分别满足不等式，且
\(\binom\alpha\beta=\alpha!/[\beta!(\alpha-\beta)!]\)。这里 \(D^0\) 为恒等操作。

\begin{definition}\label{b1:def:test-function}
设 \(\varphi:\Omega\rightarrow\mathbb K\) 光滑。如果其支集
\[
 \operatorname{supp}\varphi
 \coloneqq\overline{\{\,x\in\Omega:\varphi(x)\ne0\,\}}^{\,\Omega}
\]
为 \(\Omega\) 的紧子集，则称 \(\varphi\) 为 \(\Omega\) 上的\term[ceshi-hanshu]{测试函数}[test function]。全体测试函数构成的线性空间记作
\[
 \mathcal D(\Omega)\coloneqq C_c^\infty(\Omega).
\]
\end{definition}
\symidx{\(\mathcal D(\Omega)\)：紧支集光滑测试函数空间}

测试函数也有“试验函数”这一译名；《分析学》第2版中译本第6.2节采用后者\footnote{该译名及节次可核对\href{https://academic.hep.com.cn/br/CN/chapter/978-7-04-017381-9/chapter06}{高等教育出版社公布的第六章目录}。这里仅借目录说明名称，不以目录作为下文拓扑构造的依据。}\termidx[shiyan-hanshu]{试验函数}。

\ProseParagraphBreak

支集的紧性和它位于开集内部同样重要。例如 \(\Omega=(0,1)\) 上的常函数 \(1\) 不是测试函数，尽管它处处光滑且有界。测试函数在靠近 \(\partial\Omega\) 的某个区域内为零，因此在 \(\Omega\) 外补零后仍是 \(\mathbb R^d\) 上的光滑函数。以后对测试函数取全空间上确界或积分时，均使用这个零延拓。

\ProseParagraphBreak

这种函数并非只有零函数。令
\[
 \vartheta(t)=
 \begin{cases}e^{-1/t},&t>0,\\0,&t\leq0.\end{cases}
\]
在正半轴上，每阶导数均是 \(e^{-1/t}\) 与 \(1/t\) 的某个多项式的乘积。由 \(s^Ne^{-s}\to0\)，这些导数在 \(t\downarrow0\) 时都趋于零；逐阶检查零点的差商，得到 \(\vartheta\in C^\infty(\mathbb R)\)。于是 \(x\mapsto\vartheta(1-|x|^2)\) 是支集为闭单位球、在开单位球内严格为正的光滑函数。平移与缩放便能在 \(\Omega\) 内任一球上放置这样的函数。

\subsection{基本操作}
\label{b1:subsec:200102}

若 \(\varphi\in\mathcal D(\Omega)\)，则 \(D^\alpha\varphi\in\mathcal D(\Omega)\)，且微分不会扩大支集。若 \(a\in C^\infty(\Omega)\)，则 \(a\varphi\) 也是测试函数；即使 \(a\) 无界，其在 \(\operatorname{supp}\varphi\) 附近的每个固定阶导数仍有界。Leibniz 公式为
\begin{equation}\label{b1:eq:test-leibniz}
 D^\alpha(a\varphi)
 =\sum_{\beta\leq\alpha}\binom\alpha\beta
 D^\beta a\,D^{\alpha-\beta}\varphi.
\end{equation}
这由一阶乘积法则对 \(|\alpha|\) 归纳得到，二项系数记录各次微分落在两个因子上的分配方式。

\ProseParagraphBreak

对 \(K\Subset\Omega\) 记
\[
 \mathcal D_K=\{\,\varphi\in\mathcal D(\Omega):\operatorname{supp}\varphi\subset K\,\},
 \quad p_{K,m}(\varphi)=\max_{|\alpha|\leq m}\sup_{x\in K}|D^\alpha\varphi(x)|.
\]
上确界可等价地取遍 \(\Omega\)，空紧集上的半范数取零。因为导数在 \(K\) 外为零，\eqref{b1:eq:test-leibniz}给出
\[
 p_{K,m}(a\varphi)\leq C_{a,K,m}\cdot p_{K,m}(\varphi),
 \quad p_{K,m}(D^\alpha\varphi)\leq p_{K,m+|\alpha|}(\varphi).
\]
常数只涉及 \(a\) 在 \(K\) 上不超过 \(m\) 阶的导数。

\ProseParagraphBreak

再记 \((\tau_h\varphi)(x)=\varphi(x-h)\)。当 \(h\to0\) 且充分小时，这些平移的支集落在 \(\Omega\) 的同一个紧子集中。各阶导数的一致连续性给出 \(D^\alpha(\tau_h\varphi-\varphi)\to0\) 一致成立。沿坐标方向还可用一维积分公式写出
\[
 \frac{\varphi(x-he_i)-\varphi(x)}h
 =-\int_0^1\partial_i\varphi(x-the_i)\,\mathrm dt.
\]
对任意阶导数重复这一等式，便知差商以相同方式趋于 \(-\partial_i\varphi\)。这比只知道每个点的差商收敛更强：后面连续泛函正需要这种各阶一致控制。

\begin{lemma}\label{b1:lem:smooth-cutoff}
设 \(K\Subset\Omega\)，\(U\subset\Omega\) 为包含 \(K\) 的开集。存在 \(\chi\in\mathcal D(U)\)，满足 \(0\leq\chi\leq1\)，且 \(\chi\) 在 \(K\) 的一个邻域内恒等于 \(1\)。
\end{lemma}
\begin{proof}
由上面的 \(\vartheta\) 构造光滑阶跃
\[
 b(t)=\frac{\vartheta(2-t)}{\vartheta(2-t)+\vartheta(t-1)}.
\]
分母处处为正，且 \(b=1\) 于 \((-\infty,1]\)，\(b=0\) 于 \([2,\infty)\)。选有限个球 \(B(x_j,r_j)\) 覆盖 \(K\)，使 \(\overline B(x_j,\sqrt2r_j)\subset U\)。令 \(b_j(x)=b(|x-x_j|^2/r_j^2)\)，并取
\[
 \chi=1-\prod_j(1-b_j).
\]
每个 \(b_j\) 在其小球上等于 \(1\)，在大球外为零。因此 \(\chi\) 具有所需值域、支集和邻域恒等性质。空紧集情形取 \(\chi=0\)。
\end{proof}

光滑截断把一个局部对象变成紧支集对象，但不会在截断恒等于 \(1\) 的区域内改变它。这个局部性质比截断函数的具体公式更常用。

\subsection{单位分解}
\label{b1:subsec:200103}

\begin{definition}\label{b1:def:smooth-partition-unity}
设 \((U_\lambda)_{\lambda\in\Lambda}\) 为 \(\Omega\) 的开覆盖。如果一族非负光滑函数 \((\psi_j)\) 的支集在 \(\Omega\) 中局部有限，每个支集包含于某个 \(U_{\lambda(j)}\)，并且
\[
 \sum_j\psi_j(x)=1\quad(x\in\Omega),
\]
则称它为从属于该覆盖的\term[guanghua-danwei-fenjie]{光滑单位分解}[smooth partition of unity]。局部有限指每个点有一个邻域，只与有限个支集相交。
\end{definition}

\begin{theorem}\label{b1:thm:smooth-partition-unity}
每个欧氏开集的开覆盖都有一族可数、从属于该覆盖的光滑单位分解，而且每个分解函数均有紧支集。
\end{theorem}
\begin{proof}
取紧集穷竭
\[
 K_j=\{\,x\in\Omega:|x|\leq j,\ \operatorname{dist}(x,\mathbb R^d\setminus\Omega)\geq1/j\,\},
 \quad j\geq1,
\]
空补集的距离约定为 \(+\infty\)，并令 \(K_0=K_{-1}=\varnothing\)。这些紧集满足 \(K_j\subset\operatorname{int}K_{j+1}\)，且其内部之并为 \(\Omega\)。紧集
\(A_j=K_j\setminus\operatorname{int}K_{j-1}\) 包含在开集 \(\operatorname{int}K_{j+1}\setminus K_{j-2}\) 内。

对每个 \(A_j\)，选有限个小球覆盖它，使各小球的稍大闭球仍包含在 \(\operatorname{int}K_{j+1}\setminus K_{j-2}\) 与某个 \(U_\lambda\) 的交中。由\cref{b1:lem:smooth-cutoff}，可取在小球闭包上等于 \(1\)、支集在相应大球内的非负光滑函数 \(\eta_{j,l}\)。这族函数的支集局部有限：给定点先取其邻域包含在某个 \(\operatorname{int}K_N\) 中，所有 \(j\geq N+2\) 的支集均避开该邻域，其余层只有有限个函数。

各层小球覆盖 \(\Omega\)，故局部有限和 \(s=\sum_{j,l}\eta_{j,l}\) 光滑且处处为正。令 \(\psi_{j,l}=\eta_{j,l}/s\)，再把双指标排成一个序列。除以正光滑函数不扩大支集，所需各性质随即成立。
\end{proof}

\begin{corollary}\label{b1:cor:precompact-smooth-partition}
若 \(\Omega\subseteq\mathbb R^d\) 为非空开集，则存在可数的光滑单位分解 \((\psi_i)\) 及局部有限开覆盖 \((V_i)\)，使
\[
 \psi_i\in C_c^\infty(V_i),\qquad V_i\Subset\Omega.
\]
特别地，每个紧集只与有限多个 \(V_i\) 相交。
\end{corollary}
\begin{proof}
在上一证明中保留承载各个 \(\eta_{j,l}\) 的稍大开球，并与相应的 \(\psi_{j,l}\) 一同枚举。这些开球的闭包均包含在 \(\Omega\) 中；每层只有有限个，而足够远的层避开任意给定紧集，故它们构成局部有限开覆盖。归一化不扩大支集，所以相应的 \(\psi_i\) 属于 \(C_c^\infty(V_i)\)。
\end{proof}

若 \(\varphi\) 是测试函数，局部有限性和 \(\operatorname{supp}\varphi\) 的紧性使
\[
 \varphi=\sum_j\psi_j\varphi
\]
实际上只有有限个非零项。于是可以先在每个覆盖开集上验证一个关于测试函数的线性恒等式，再把有限项相加恢复全局恒等式。局部有限的要求不能换成“可数”二字；一般可数和既未必光滑，也未必允许逐项操作。

\subsection{测试函数空间拓扑}
\label{b1:subsec:200104}

先在 \(\mathcal D_K\) 上使用半范数 \(p_{K,m}\) 所给的拓扑。其收敛就是各阶导数一致收敛，可由度量
\[
 \mathrm{d}_K(\varphi,\psi)=\sum_{m=0}^\infty2^{-m-1}\min\{1,p_{K,m}(\varphi-\psi)\}
\]
描述。这个空间完备。事实上，Cauchy 列的每个导数都有一致极限 \(g_\alpha\)。在包含于 \(\Omega\) 的小球中，将普通积分恒等式
\[
 D^\alpha\varphi_j(x+he_i)-D^\alpha\varphi_j(x)
 =\int_0^h D^{\alpha+e_i}\varphi_j(x+te_i)\,\mathrm dt
\]
传到极限，得到 \(\partial_i g_\alpha=g_{\alpha+e_i}\)。故 \(g_0\) 光滑，其支集仍在 \(K\)，且原列在每个半范数中收敛到它。

\ProseParagraphBreak

整个 \(\mathcal D(\Omega)\) 是这些固定支集空间的并，却不能直接赋予 \(C^\infty(\Omega)\) 的局部一致拓扑。例如，在 \(\Omega=\mathbb R^d\) 中，一个固定小凸起向无穷远平移，在每个紧集上终将为零，但其积分始终不变。若这样的列在测试空间中趋于零，积分泛函就会失去连续性。

\ProseParagraphBreak

以下从半范数族构造完整局部凸拓扑的部分属于技术性补充；第一次阅读可以先接受随后\cref{b1:prop:test-sequence-convergence}与\cref{b1:prop:test-functional-continuity}给出的序列判据和有限阶估计，再回读这一构造。为避免把一个收敛列的判据误当成对全部开集的定义，设 \(\mathscr P\) 为 \(\mathcal D(\Omega)\) 上所有这样的半范数 \(q\)：对每个 \(K\Subset\Omega\)，存在 \(C>0\) 与整数 \(m\geq0\)，使
\[
 q(\varphi)\leq C\cdot p_{K,m}(\varphi)\quad(\varphi\in\mathcal D_K).
\]
以有限个 \(q\in\mathscr P\) 的小于正数的集合之交作为零点邻域基，定义 \(\mathcal D(\Omega)\) 的拓扑。半范数的不等式保证加法、数乘连续，而且每个 \(\mathcal D_K\) 的包含映射连续。由于这些邻域都是凸的，这是一种局部凸拓扑；它也是使各包含映射连续的最强局部凸拓扑。这最后一句只须注意：任何另一个由半范数给出的此类拓扑，其定义半范数在每个 \(\mathcal D_K\) 上连续，因而都已收入 \(\mathscr P\)。

\begin{proposition}\label{b1:prop:test-sequence-convergence}
序列 \(\varphi_j\) 在 \(\mathcal D(\Omega)\) 中趋于零，当且仅当存在同一个 \(K\Subset\Omega\) 包含所有 \(\operatorname{supp}\varphi_j\)，并且对每个 \(m\geq0\)，有 \(p_{K,m}(\varphi_j)\to0\)。
\end{proposition}
\begin{proof}
若具有右侧性质，则每个 \(q\in\mathscr P\) 的固定紧集估计给出 \(q(\varphi_j)\to0\)，所以在所定义拓扑中收敛。

反之，假定 \(\varphi_j\to0\)。全空间上各阶导数的上确界半范数属于 \(\mathscr P\)，故这些量趋于零。还须证明支集可统一限制。若不能，对上面紧集穷竭可选一个子列 \(\varphi_{j_l}\) 及点 \(x_l\notin K_l\)，使 \(\varphi_{j_l}(x_l)\ne0\)。每个有限的前缀有共同紧支集，因此指标可递增选取。点族 \((x_l)\) 在 \(\Omega\) 中局部有限。于是
\[
 q(\varphi)=\sum_l\frac{|\varphi(x_l)|}{|\varphi_{j_l}(x_l)|}
\]
是良定义的半范数：每个紧支集只遇到有限个 \(x_l\)，在固定 \(\mathcal D_K\) 上有 \(q\leq C_K\cdot p_{K,0}\)。故 \(q\in\mathscr P\)，但 \(q(\varphi_{j_l})\geq1\)，与收敛矛盾。这证明共同紧支集存在。
\end{proof}

\begin{proposition}\label{b1:prop:test-functional-continuity}
设 \(T:\mathcal D(\Omega)\rightarrow\mathbb K\) 线性。下列三个条件等价：
\begin{equivalences}
\item\label{b1:item:test-functional-topological}泛函 \(T\) 在上述拓扑中连续；
\item\label{b1:item:test-functional-sequential}每当 \(\varphi_j\to0\) 于 \(\mathcal D(\Omega)\)，都有 \(T\varphi_j\to0\)；
\item\label{b1:item:test-functional-order}对每个 \(K\Subset\Omega\)，存在 \(C_K>0\)、整数 \(m_K\geq0\)，使
\[
 |T\varphi|\leq C_K\cdot p_{K,m_K}(\varphi)\quad(\varphi\in\mathcal D_K).
\]
\end{equivalences}
\end{proposition}
\begin{proof}
连续性蕴涵序列连续性。若第三项对某个 \(K\) 失败，则对每个 \(j\geq1\) 可取 \(u_j\in\mathcal D_K\)，满足 \(|Tu_j|>j\cdot p_{K,j}(u_j)\)。令 \(v_j=u_j/|Tu_j|\)，则对每个固定的 \(m\)，在 \(j\geq m\) 时有 \(p_{K,m}(v_j)<1/j\)，但 \(|Tv_j|=1\)。这与第二项矛盾。最后，第三项说明半范数 \(\varphi\mapsto|T\varphi|\) 属于 \(\mathscr P\)，从而 \(T\) 连续。
\end{proof}

这个有限阶估计是后面使用测试函数拓扑的主要接口。微分和乘以固定光滑函数都保持测试列收敛；在线性泛函上，这足以证明转移过去的相应操作仍连续。紧集改变时允许估计中的阶数和常数一起改变，没有理由要求一个无界开集上的所有位置受同一有限阶控制。
